%% LaTeX-Beamer template for KIT design
%% by Erik Burger, Christian Hammer
%% title picture by Klaus Krogmann
%%
%% version 2.0
%%
%% mostly compatible to KIT corporate design v2.0
%% http://intranet.kit.edu/gestaltungsrichtlinien.php
%%
%% Problems, bugs and comments to
%% burger@kit.edu

%\documentclass[18pt]{beamer}
\documentclass[18pt, aspectratio=169]{beamer}


\usepackage{xcolor}
\usepackage{hyperref}

\usepackage{pgfpages}
%\setbeameroption{show notes on second screen = right}
%\documentclass[18pt,draft]{beamer}

\usetheme{kit2}
\setbeamercovered{invisible}
\setbeamertemplate{navigation symbols}{}


%%%% Zeichencodierung & Schrift
\usepackage[ansinew]{inputenc}
\usepackage[ngerman]{babel}
\usepackage{amssymb}
\let\raggedright\relax
\usepackage{ragged2e}
%\usepackage{color}


%%%% Einbinden von Graphiken
%\usepackage{graphicx}
\usepackage{graphics}
\usepackage{graphicx}
\usepackage{pgf,tikz}


%%%% Bibliographie und Zitationen
\usepackage[round]{natbib}
\bibliographystyle{nachvorj}


%%%% Tabellen
\usepackage{multirow,booktabs,verbatim }
\setbeamertemplate{caption}[numbered]

%\graphicspath{{U:/Lehre/MethodenII/Grafiken/}}
%\graphicspath{{D:/uni/Lehre/MethodenII/Grafiken/}}
\graphicspath{{/Dropbox/Graphs/}}



\hypersetup{pdfpagemode=FullScreen}


\setbeamerfont{block title}{size=\normalsize}

\title{Datenauswertung}
\subtitle{Multivariate Regression - Komplexere F�lle}
\author{Andreas Haupt \\ andreas.haupt@kit.edu }
\institute{Karlsruher Institut f�r Technologie}
\begin{document}

% change the following line to "ngerman" for German style date and logos
\selectlanguage{ngerman}

%title page
\begin{frame}
\titlepage
\end{frame}



%%%%%%%%%%%%%%%% Interaktion %%%%%%%%%%%%%%%%%%%%%%%%%%%%%%%

\section{Interaktion}
\subsection*{}
\begin{frame}
\frametitle{Interaktion}
\begin{figure}
\centering
\begin{tikzpicture}
\draw[line width=1pt, ->] (0,3) -- (0,0);
\draw[line width=1pt, ->] (-2.5,-0.25) -- (2.5,-0.25);
\node at (0,3.25) {\large Z};
\node at (0,-0.25) {\large \textit{I}};
\node at (-2.75,-0.25) {\large X};
\node at (2.75,-0.25)  {\large Y};
\end{tikzpicture}
\end{figure}
Eine Interaktion liegt genau dann vor, wenn die St�rke des Effekts von X auf Y systematisch von einer dritten Eigenschaft Z abh�ngig ist. Dabei ist es unerheblich, ob die Eigenschaft Z \textit{zus�tzlich} \textit{auch} einen Einfluss auf X und Y aus�bt. 
\end{frame}


\frame{
\frametitle{Interaktion}

\begin{center}
\resizebox{.7\linewidth}{!}{
  \includegraphics{lehre-ss20-frau-bildung-lesen}
  }
\end{center}

}






\begin{frame}[fragile]
\frametitle{Interaktion}

Wenn wir mittlere Bildung ($B$) als Referenzkategorie f�r die Bildung nutzen und das Geschlecht ($F$) so codieren, dass M�nner eine 0 erhalten und Frauen eine 1, k�nnten wir versucht sein, folgende Regression zu sch�tzen: 

\[ \hat{y}_i = \hat{\alpha} + \hat{\beta_1}F_i + \hat{\beta_2}B_{i_{niedrig}} + \hat{\beta_3}B_{i_{hoch}}  \] 

Das w�rde uns einen Geschlechtereinfluss $\hat{\beta_1}F_i$ als auch einen vom Geschlecht \textbf{unabh�ngigen} Einfluss von Bildung geben. Wir m�chten aber einen \textbf{geschlechtsspezifischen} Einfluss berechnen. 

\end{frame}

\begin{frame}[fragile]
\frametitle{Fall 5: Interaktion}

Wir m�ssen daher die Terme f�r Bildung f�r beide Geschlechter seperat berechnen. Das kann man auf folgende Art erreichen: 

\[ \hat{y}_i = \hat{\alpha} + \hat{\beta_1}F_i + \hat{\beta_2}B_{i_{niedrig}} + \hat{\beta_3}B_{i_{hoch}} + \hat{\beta_4}B_{i_{niedrig}}\cdot F_i + \hat{\beta_5}B_{i_{hoch}}\cdot F_i  \] 

Damit k�nnen wir alle sechs Gruppen abbilden: 

\begin{center}
{{ \small 
\begin{tabular}{lcccr} \toprule
                    & Geschlecht & niedrige Bildung & hohe Bildung & Gruppe\\  \midrule
$\hat{\alpha}$                          & 0 & 0 & 0 & M�nner, mittel \\
$\hat{\beta_1}F_i$                      & 1 & 0 & 0 & Frauen, mittel \\
$\hat{\beta_2}B_{i_{niedrig}}$          & 0 & 1 & 0 & M�nner, niedrig \\
$\hat{\beta_3}B_{i_{hoch}}$             & 0 & 0 & 1 & M�nner, hoch \\
$\hat{\beta_4}B_{i_{niedrig}} \cdot F_i$       & 1 & 1 & 0 & Frauen, niedrig \\
$\hat{\beta_5}B_{i_{hoch}} \cdot F_i$          & 1 & 0 & 1 & Frauen, hoch\\ \bottomrule
\end{tabular}
}}
\end{center}
\end{frame}

\begin{frame}[fragile]
\frametitle{Fall 5: Interaktion}
\begin{center}
  \begin{minipage}[c]{.3\linewidth} % [b] => Ausrichtung an \caption
\begin{center}
\resizebox*{1\linewidth}{!}{
\input{/Dropbox/Tabellen/lehre-ss20-reg-bild-frau.tex}
}
\end{center}
\vfill
\quad \\[8em]
  \end{minipage}
  \hspace{.15\linewidth}% Abstand zwischen Bilder
  \begin{minipage}[b]{.5\linewidth} % [b] => Ausrichtung an \caption 
  \begin{center}
\resizebox{1\linewidth}{!}{
  \includegraphics{lehreSS20-lesereg-interaktV2}
}
\end{center}
\end{minipage}
\end{center}
\
\end{frame}

%%%%%%%%%%%%%%%% Collider %%%%%%%%%%%%%%%%%%%%%%%%%%%%%%%

\section{Collider}
\subsection*{}

\begin{frame}
\frametitle{The Dark Triangle: Collider}
\begin{figure}
\centering
\begin{tikzpicture}
\draw[line width=1pt, ->] (-2.65,0) -- (-0.25,3);
\draw[line width=1pt, <-] (0.25,3) -- (2.65,0);
\draw[line width=1pt, ->] (-2.5,-0.25) -- (2.5,-0.25);
\node at (0,3.25) {\large Z};
\node at (-2.75,-0.25) {\large X};
\node at (2.75,-0.25)  {\large Y};
\end{tikzpicture}
\end{figure}

{{Ein Collider liegt dann vor, wenn ein Treatment X einen Einfluss auf das Outcome Y hat, aber parallel die Verteilung einer dritten Eigenschaft Z vom Outcome bedingt wird.  In einem solchen Fall f�hrt die Isolation des Einflusses von X $\rightarrow$ Y mit Hilfe von Z zu v�llig falschen Ergebnissen.}}

\end{frame}


\begin{frame}{Medienkonsum und �bergewicht (again)}
\begin{columns}

\small

\begin{column}{.4\linewidth}
\begin{center}
 \resizebox{1\linewidth}{!}{
 \includegraphics{../Graphs/screentime-obesity} 
 }
\end{center}
{{ \scriptsize Source: \href{https://pubmed.ncbi.nlm.nih.gov/29967527/}{Williams et al. (2018)}}}

\end{column}

\begin{column}{.6\linewidth}
\begin{itemize}
	\item Selbstverletzungstendenzen bilden in Modell d ein Collider mit �bergewicht und Medienkonsum. 
	\item Laut diesem Modell neigen Personen mit �bergewicht st�rker zu Selbstverletzungen und sind daher in der Gruppe der ``Selbstverletztenden'' �berrepr�sentiert. 
	\item Das Modell nimmt au�erdem einen direkten Einfluss des Medienkonsums auf die Selbstverletzungstendenz an. 
	\item Wenn wir den Einfluss von Medienkonsum auf �bergewicht isolieren wollen, d�rfen wir nicht f�r Selbstverletzungstendenzen kontrollieren, weil die Unterscheidung nach diesen Tendenzen \textit{eine implizite, erneute} Unterscheidung nach �bergewicht ist.
\end{itemize}
\end{column}

\end{columns}
\end{frame}


\begin{frame}{F�rdern Bibliotheksbesuche das Leseinteresse?}
\begin{columns}

\small

\begin{column}{.4\linewidth}
\begin{center}
 \resizebox{1\linewidth}{!}{
 \includegraphics{../Graphs/Leseleistung-Kollider} 
 }
\end{center}
\end{column}

\begin{column}{.6\linewidth}
\begin{itemize}
	\item Im linken Modell gehen Sch�ler:innen mit h�herem Leseinteresse mit h�herer Wahrscheinlichkeit in die Bibliothek. 
	\item Frauen gehen mit h�herer Wahrscheinlichkeit in die Bibliothek. 
	\item In diesem Modell erh�ht ein Bibliotheksbesuch das Leseinteresse nicht!
\end{itemize}
\end{column}

\end{columns}
\end{frame}


\begin{frame}
\frametitle{The Dark Triangle: Collider}
\begin{center}
  \begin{minipage}[b]{.35\linewidth} % [b] => Ausrichtung an \caption
\begin{center}
\resizebox{1\linewidth}{!}{
\includegraphics{lehre-ss20-coll-lesint-frau}
}
\end{center}
  \end{minipage}
  \hspace{.05\linewidth}% Abstand zwischen Bilder
  \begin{minipage}[b]{.35\linewidth} % [b] => Ausrichtung an \caption
\begin{center}
\resizebox{1\linewidth}{!}{
\includegraphics{lehre-collider-bib-lesint}
}
\end{center}
  \end{minipage}
  \hspace{.05\linewidth}% Abstand zwischen Bilder
  \begin{minipage}[b]{.35\linewidth} % [b] => Ausrichtung an \caption
\begin{center}
\resizebox{1\linewidth}{!}{
\includegraphics{lehre-ss20-coll-lesint-frau-bib}
}
\end{center}
  \end{minipage}
\end{center}
\end{frame}

\begin{frame}
\frametitle{The Dark Triangle: Collider}

\begin{tabular}{l*{6}{c}}
\toprule
                    &\multicolumn{1}{c}{(1)}&\multicolumn{1}{c}{(2)}&\multicolumn{1}{c}{(3)}&\multicolumn{1}{c}{(4)}&\multicolumn{1}{c}{(5)}&\multicolumn{1}{c}{(6)}\\
                    &\multicolumn{3}{c}{Leseinteresse}&\multicolumn{3}{c}{Bibliotheksnutzung}\\
\midrule
Frau                &        20.0\sym{***}&                     &        14.1\sym{***}&        20.0\sym{***}&                     &        13.6\sym{***}\\
                    &      (0.51)         &                     &      (0.50)         &      (0.53)         &                     &      (0.52)         \\
Bibliotheksnutzung  &                     &        32.7\sym{***}&        29.5\sym{***}&                     &                     &                     \\
                    &                     &      (0.52)         &      (0.53)         &                     &                     &                     \\
Leseinteresse       &                     &                     &                     &                     &         0.4\sym{***}&         0.3\sym{***}\\
                    &                     &                     &                     &                     &      (0.01)         &      (0.01)         \\
Konstante           &       120.9\sym{***}&       108.9\sym{***}&       103.9\sym{***}&        57.4\sym{***}&        21.3\sym{***}&        18.8\sym{***}\\
                    &      (0.36)         &      (0.43)         &      (0.46)         &      (0.38)         &      (0.78)         &      (0.78)         \\
\midrule
N        &       30000         &       30000         &       30000         &       30000         &       30000         &       30000         \\
\bottomrule
\multicolumn{6}{l}{\footnotesize Standardfehler in Klammern}\\
\multicolumn{6}{l}{\footnotesize \sym{*} \(p<0.05\), \sym{**} \(p<0.01\), \sym{***} \(p<0.001\)}\\
\end{tabular}

\end{frame}

\begin{frame}{Collider und Stichprobenselektion}
\begin{columns}
\small

\begin{column}{.5\linewidth}
\begin{center}
 \resizebox{1\linewidth}{!}{
 \includegraphics{../Graphs/dag-collider-depression} 
 }
\end{center}
{{\scriptsize Source: \href{https://www.aerzteblatt.de/archiv/223244/Collider-Bias-in-Beobachtungsstudien-Konsequenzen-fuer-die-medizinische-Forschung}{T�nnies et al. (2022)}}}
\end{column}

\begin{column}{.5\linewidth}
\begin{itemize}
	\item Die Probleme ``Kontrolle mit einem Collider'' und Stichprobenselektion k�nnen sich �berlagern. 
	\item Prozesse, die mit einer Stichprobenselektion einhergehen, k�nnen auch mit dem Treatment, einer Kontrollvariable und/oder dem Outcome korreliert sein.
	\item In der Grafik bildet einen Fall ab, in der Personen mit geringeren depressiven Zust�nden und besserer Gesundheit eher an einer Studie teilnehmen. Dadurch ist die analytische Stichprobe im Vergleich zur Population selektiv. Die Eigenschaft ``Studienteilnahme'' ist daher ein Collider.
\end{itemize}
\end{column}

\end{columns}
\end{frame}



\section{Gute und schlechte Modelle}
\subsection*{}

\begin{frame}{Effekte auf einen Blick}
\begin{columns}

\begin{column}{.5\linewidth}
\begin{center}
 \resizebox{1\linewidth}{!}{
 \includegraphics{../Graphs/rohrer-dag-explainer} 
 }
\end{center}
{{ \scriptsize Quelle: \href{https://compass.onlinelibrary.wiley.com/doi/pdf/10.1111/spc3.12948}{Rohrer (2024, Abbildung 1)}}}
\end{column}

\begin{column}{.5\linewidth}
\begin{itemize}
	\item Soziale Prozesse k�nnen weitaus komplexer sein, als Dreiecksbeziehungen. 
	\item Typischerweise k�nnen Sie aber jedes einzelne Element in komplexen Beziehungen in Form von Dreiecksbeziehungen darstellen. 
	\item Reduzieren Sie daher Komplexit�t, indem Sie stets die Ihr Ziel im Auge behalten: Sie wollen gute Vergleiche. 
\end{itemize}
\end{column}

\end{columns}
\end{frame}

\begin{frame}
\frametitle{Leseempfehlungen}

�ber die graphische Darstellung von Kausaltheorien: \href{https://cran.r-project.org/web/packages/ggdag/vignettes/intro-to-dags.html}{Intro to DAGs}

\vfill

Ein kleiner ``Crash Course'' �ber gute und schlechte Wege Einfl�sse zu isolieren: \href{http://causality.cs.ucla.edu/blog/index.php/2019/08/14/a-crash-course-in-good-and-bad-control/}{A Crash Course in Good and Bad Control}

\end{frame}

\end{document}
